\exerciseeditionsection{Solutions to Wald's Problems}
\ExerciseEditorialNote

\begin{enumerate}[leftmargin=2.8em,itemsep=1.2em]
\WaldSolution{1}
\begin{enumerate}[label=\alph*),leftmargin=2em]
\item Differentiating \(E=-u_a\xi^a\) and using the Killing equation gives
\[
 \frac{\dd E}{\dd\tau}=-a_a\xi^a.
\]
Decompose \(\xi^a=Eu^a+s^a\), with \(s^au_a=0\).  Since
\(s_as^a=E^2+\xi_a\xi^a<E^2\), and \(a^au_a=0\), Cauchy--Schwarz on the
spatial rest space gives
\[
 \left|\frac{\dd E}{\dd\tau}\right|=|a_as^a|
 \le a|s|\le aE.
\]

\item Division by \(E>0\) yields
\( |\dd\ln E/\dd\tau|\le a\).  Thus between any two points of the curve,
\[
 E_1e^{-\int a\dd\tau}\le E_2
 \le E_1e^{\int a\dd\tau}.
\]
Finite integrated acceleration bounds \(E\).  The reverse Lorentzian
Cauchy--Schwarz inequality gives
\(E=-u\mathbin\cdot\xi\ge\sqrt{-\xi^2}\), so
\(-\xi^2\le E^2\) is bounded along the curve.
\end{enumerate}

\WaldSolution{2}
On \(x=\pi/2\) or \(x=3\pi/2\), the vector \(k^a=(\dd y/\dd\lambda)
(\partial/\partial y)^a\) is null.  Substitution in the two coordinate
geodesic equations shows that the \(x\)-equation is satisfied and the
remaining equation is
\[
 \frac{\dd^2y}{\dd\lambda^2}
 +\frac12\left(\frac{\dd y}{\dd\lambda}\right)^2=0
\]
after choosing the indicated orientation.  Integration gives
\[
 \frac12\lambda\frac{\dd y}{\dd\lambda}=1,
 \qquad y=2\ln|\lambda|+\text{constant}.
\]
As \(\lambda\to0^+\), \(y\to-\infty\), which represents infinitely many
negative windings around the periodic \(y\)-coordinate in finite affine
parameter.  A change \(y\mapsto y-2\pi\) rescales \(\lambda\) by
\(e^{-\pi}\), so the affinely normalized tangent returns larger by
\(e^\pi\) after one negative cycle.

\WaldSolution{3}
\begin{enumerate}[label=(\roman*),leftmargin=2em]
\item Let \(p_n\to p\) in the compact horizon.  Lower semicontinuity of
\(T\) follows by slightly varying the maximizing timelike geodesic from
\(p\) to \(S\).  For upper semicontinuity, compactness and the limit-curve
argument give a convergent subsequence of maximizing geodesics
\(\gamma_{p_n}\); its limit joins \(p\) to \(S\), and upper semicontinuity
of Lorentzian length bounds the limiting lengths by \(T(p)\).  Hence
\(T(p_n)\to T(p)\).  Compactness of \(H^-(S)\) then gives a point where
\(T\) is minimal.

\item Let \(q\) lie to the future of \(p\) on a null generator.  Concatenate
the null segment from \(p\) to \(q\) with a maximizing timelike geodesic
from \(q\) to \(S\).  The corner can be pushed to a timelike curve from
\(p\) to \(S\) whose length exceeds the timelike portion, so
\(T(p)>T(q)\).  A future continuation of the generator from a minimum point
would therefore have a smaller value of \(T\), a contradiction.  The strong
causality hypothesis is unnecessary.
\end{enumerate}
\end{enumerate}
